\section{Axis-Based Representation}
\label{sec:axis-model}
Within the defined scope (Section~\ref{sec:formal-scope}), each computation
component is determined by its input dimension connections and output axis order.

\subsection{Logical Axes}
\label{sec:logical-axes}
A \emph{dimension} is a physical tensor dimension, selected by a zero-based
position in an \code{axis} attribute. A \emph{logical axis} is an equivalence
class of input dimensions that use the same coordinate. For example, in
\code{einsum('ik,kj->ij', A, B)}~\cite{numpy-einsum}, the dimensions marked
\code{k} share a coordinate: the second dimension of $A$ and the first dimension of
$B$. Since \code{k} is absent from the output \code{ij}, it is summed out.
We record this relation using input \emph{ports}. Port $(X,d)$ is dimension
$d$ of input $X$, with positions starting at zero. For a component with input
set $\mathcal X$, let
\begin{equation}
 P=\{(X,d)\mid X\in\mathcal X,\ 0\le d<\operatorname{rank}(X)\}.
\end{equation}
We write $u\sim v$ when ports $u$ and $v$ use the same coordinate.
Each logical axis is one equivalence class. For the LHS expression $L$, the
set of all logical axes is
\begin{equation}
 \Axes(L)=P/{\sim}.
\end{equation}
Here $P/{\sim}$ groups connected ports into one class each, including
classes removed by summation.
Equal extents do not establish axis identity. We also retain the classes of
axes removed by summation: they still record how input coordinates connect.

\subsection{Symbolic Axis Groups}
\label{sec:symbolic-axis-groups}
A \emph{symbolic axis group} represents an ordered sequence of distinct
logical axes whose length need not be known at compile time. We write groups
as $\vec I$, $\vec J$, or $\vec K$. For example,
$\vec K=(k_1,\ldots,k_r)$ represents $r$ axes, where $r$ may vary across
matches. This allows one symbolic description to cover tensors of different
ranks.

Group order matters because operation attributes refer to dimension
positions. We write $\vec I\Vert\vec K$ for concatenation and
$\operatorname{set}(\vec I)$ for the corresponding unordered set. When a
group appears in both operands, corresponding entries denote the same
logical axes. For example, an unbatched contraction with operand orders
$\vec I\Vert\vec K$ and $\vec K\Vert\vec J$ sums over $\vec K$ and
produces output order $\vec I\Vert\vec J$.

\subsection{Axis Representation}
\label{sec:axis-representation}
For an expression $T$ with assigned attributes, define
\begin{equation}
 N(T)=(P,\sim,F),
 \label{eq:axis-representation}
\end{equation}
where $P$ is the set of input ports, $\sim$ records their coordinate
connections, and $F$ lists the output classes in dimension order. This representation
specifies the value at each output coordinate: use the same index at
connected input ports, multiply the selected input entries, and sum over
all classes absent from $F$. Classes in $F$ remain as output coordinates.

For matrix multiplication $T=AB$, contraction identifies
$(A,1)\sim(B,0)$, giving the classes:
\begin{equation}
 i=\{(A,0)\},\qquad k=\{(A,1),(B,0)\},\qquad j=\{(B,1)\}.
\end{equation}
Thus
\begin{equation}
 N(T)=\bigl(\{(A,0),(A,1),(B,0),(B,1)\},\sim_{AB},(i,j)\bigr).
\end{equation}
The class $k$ is absent from the output, giving $T_{ij}=\sum_k A_{ik}B_{kj}$.

\subsection{Equivalence of Axis Representations}
\begin{lemma}[Axis exactness]
\label{lem:axis-exactness}
For two legal components $L,R$ over the same inputs at fixed match metadata,
with every input dimension having extent at least two,
\begin{equation}
 \forall X,\quad L(X)=R(X)\quad\iff\quad N(L)=N(R).
 \label{eq:axis-exactness}
\end{equation}
\end{lemma}
That is, within the defined scope, two computations produce the same output
shape and values for every input exactly when their input-dimension connections
and output-axis order agree. The names assigned to logical axes do not matter.

For TEPL, this gives attribute synthesis a concrete target: choose legal RHS
axis fields that reproduce the LHS axis representation. The lemma establishes
that matching this representation is both necessary and sufficient for
equivalence, providing the basis for the synthesis algorithm in
Section~\ref{sec:mathematical-synthesis}. Appendix~\ref{app:axis-exactness}
gives the proof and a concrete example.
